\chapter{Conclusion and Future Scope}
\label{ch:conclusion}

\section{Summary of the Study}
This dissertation has presented the design, implementation and
evaluation of a comprehensive Q\&A based interactive storytelling
system powered by Large Language Models. The system positions users
as active co-creators rather than passive consumers of narrative
content. Through a structured 17-question elicitation protocol
across four narrative dimensions (setting, character, theme, plot),
a seven-module backend pipeline, and a provider-agnostic LLM layer,
the system transforms user preferences into coherent, personalised,
multi-chapter stories.

Beyond text, the system delivers a multimodal reader experience:
per-segment scene illustrations are synthesised on demand via
Pollinations.ai using a sentence-scored prompt extractor and a
serial-throttled request queue; spoken narration is produced through
the W3C Web Speech Synthesis interface; and genre-conditioned ambient
music is streamed from a royalty-free CDN. Reader analytics surface
through a D3-based Story Map and a Character Relationship Graph
that openly disclaim their co-appearance basis. An end-user
authoring layer admits free-form custom answers and entirely
user-added Q\&A pairs, routed through a dedicated
\texttt{custom\_elements} context channel into the LLM prompt.

\section{Main Findings}
\begin{itemize}
    \item Across 50 generated stories and 15 user-study participants,
    the system achieved mean Likert scores of 4.05 / 5 (coherence),
    4.12 / 5 (creativity), 4.35 / 5 (engagement) and 3.75 / 5
    (consistency).
    \item Groq's free LLaMA-3.1-8B cloud inference delivered a 1.8\,s
    mean initial-segment latency, dramatically faster than local
    Ollama (12.4\,s) and competitive with paid OpenAI GPT-3.5-turbo
    (2.6\,s).
    \item The 12-question budget emerged as the recommended default:
    it pushes personalisation past 4.3 / 5 while keeping elicitation
    time under three minutes.
    \item The serial throttling queue reduced image-API failure
    probability from $\approx$70\% to $<$2\% under 5$\times$ parallel
    load.
    \item Filtering and a min-2 occurrence threshold raised
    character-extraction $F_1$ from 0.71 to 0.80 while halving the
    false-positive rate.
\end{itemize}

\section{Contributions to Knowledge}
\begin{enumerate}
    \item \textbf{Novel Q\&A framework} for interactive
    storytelling that enables real-time personalisation through
    conversational interaction across four canonical narrative
    dimensions.
    \item \textbf{Custom-elements context channel} --- a deployable
    architectural pattern that lets end users inject unstructured
    Q\&A pairs into the LLM prompt without code or schema changes;
    generalises beyond storytelling to any LLM-driven application
    requiring extensible personalisation.
    \item \textbf{Sentence-scored visual prompt extractor} that
    mitigates dialogue and inner-thought contamination when feeding
    story text into a text-to-image model.
    \item \textbf{Cross-modal serial-throttling queue} that prevents
    Cloudflare-fronted free image services from rate-limiting
    bursty generation; image-synthesis success rate exceeds 95\%
    after retries.
    \item \textbf{Deliberately disclaimed} co-appearance reader graph,
    establishing a pattern of \emph{honest reporting} for
    visualisations whose underlying inference is statistical rather
    than semantic.
\end{enumerate}

\section{Practical Implications}
The dissertation has practical implications for several stakeholder
groups:

\begin{itemize}
    \item \textbf{Educators} can use the system to create
    personalised story-based tutorials that adapt to learner
    understanding and interest.
    \item \textbf{Game designers} can adopt the Q\&A elicitation
    pattern to capture rich player preferences without sacrificing
    narrative control.
    \item \textbf{Therapists} can pilot the system as a narrative-
    therapy adjunct that respects user agency.
    \item \textbf{Content creators} can prototype scalable
    personalised storytelling for marketing or children's literature
    at near-zero infrastructure cost.
    \item \textbf{Researchers} benefit from an open-source baseline
    \cite{jha2026repo} for future studies on interactive narrative.
\end{itemize}

\section{Limitations of the Study}
Several limitations bound the present work:

\begin{itemize}
    \item Keyword-based context extraction misses paraphrastic user
    inputs; a fine-tuned BERT/DistilBERT \cite{devlin2019bert,sanh2019distilbert}
    or LLM zero-shot classifier would improve recall.
    \item In-memory session storage limits horizontal scalability;
    Redis-backed sessions are required for distributed deployment.
    \item The consistency checker is pattern-based and lacks
    semantic grounding; cross-segment NLI \cite{williams2018mnli}
    is a natural upgrade.
    \item Image quality inherits the ceiling of Pollinations.ai
    \cite{pollinations2024api}; switching to local Stable
    Diffusion~XL \cite{podell2024sdxl} or FLUX.1 \cite{bfl2024flux}
    would improve consistency at the cost of GPU requirement.
    \item Ambient music is hot-linked rather than locally hosted,
    creating a single-point dependency on SoundHelix's availability.
    \item The character relationship graph is openly disclaimed as
    co-appearance --- a genuine relationship-extraction module
    would require per-segment NLU inference.
    \item User-study sample size ($n=15$) is modest by HCI standards
    and skews young (ages 22--35); a larger and more diverse cohort
    would strengthen external validity.
\end{itemize}

\section{Future Scope}
Building on the foundation established by this dissertation, several
directions for future research emerge:

\begin{itemize}
    \item \textbf{Semantic context extraction} via DistilBERT or
    zero-shot LLM classification, replacing keyword dictionaries.
    \item \textbf{Sentiment-aware relationship-extraction graph}
    using small NLU models per segment.
    \item \textbf{On-device image generation} with local Stable
    Diffusion~XL or FLUX.1 for offline operation and per-character
    visual consistency.
    \item \textbf{Neural TTS} (e.g.\ XTTS \cite{casanova2024xtts})
    with cloned voices for character-specific narration.
    \item \textbf{Music conditioned by current segment mood}
    (per-segment crossfade) rather than per-genre static loops.
    \item \textbf{Collaborative multi-user sessions} with merged
    Q\&A elicitation and turn-taking.
    \item \textbf{Containerisation} (Docker) plus Redis sessions
    for distributed deployment.
    \item \textbf{Automated evaluation} using BLEU, chrF, MAUVE and
    narrative-arc detection \cite{lin2024evalnarrate}.
    \item \textbf{Fine-tuned LoRA adapters} per genre on curated
    open story corpora.
\end{itemize}

\section{Final Thoughts}
The remarkable progress of LLMs has put high-quality text generation
within reach of every developer, but \emph{turning text generation
into a great storytelling experience} remains a design challenge,
not just a capability challenge. This dissertation has argued and
demonstrated that the missing ingredient is a clean separation
between \emph{what the system asks the user} and \emph{what it does
with the answers}. By making the elicitation protocol structured but
extensible, the LLM provider pluggable and the modalities
orchestrated rather than bolted on, the system delivers a
storytelling experience that is at once accessible to first-time
users and extensible by curious ones. The code is open-source
\cite{jha2026repo}; the framework is ready for adoption.
